%=====================================================================
\section{Résultats}

\subsection{La référence optique}

\paragraph{Le guide appliqué tel quel sur Sentinel-2.} Dans la forêt de G23, les seuils du guide valent
$m=-0{,}040$ et $\sigma=0{,}058$ ($\Delta$NDVI). La méthode produit 3\,349 objets de 1~ha ou plus (13\,600~ha),
soit \textbf{8,7~\% de la forêt en deux mois et demi}, alors que le taux de coupe rase est de l'ordre de
1~\% par an. Les détections sont surtout des feuillus, qui jaunissent dès septembre
(\cref{fig-s2-overview}). Les seuils $m-k\sigma$ supposent une forêt stable entre les deux dates ; ce
n'est vrai que pour deux images d'été à un an d'écart, pas pour un suivi intra-saison.

\begin{figure}[H]
\centering
\includegraphics[width=0.9\linewidth]{s2_guide_overview.png}
\caption{Méthode du guide sur Sentinel-2, toute la G23 (degrés de certitude, juillet $\to$ septembre 2026).}
\label{fig-s2-overview}
\end{figure}

\paragraph{Référence qualifiée et photo-interprétation.} En ne gardant que le passage d'un état boisé au
sol nu (degré 3, NDVI avant $\geq 0{,}70$, après $\leq 0{,}65$, tamisage à 1~ha), on obtient 281 coupes
(613~ha). La photo-interprétation d'un échantillon de 30 coupes donne une précision de 59~\% (coupes
certaines) à 86~\% (avec les probables), sur 22 objets interprétables ; toutes les erreurs sont en
feuillus. Dans la fenêtre radar (\cref{fig-window}), les 60 coupes qualifiées ont toutes été examinées :
seules \textbf{10 montrent un sol nu certain} (1 en feuillus, 6 en résineux, 3 en mixtes), 9 sont
probables et 41 ne montrent pas de sol nu (sénescence, ombre, éclaircie). Le test \enquote{forêt avant}
(SWIR de juillet) écarte trois coupes déjà ouvertes en juillet. La référence principale compte donc
\textbf{7 coupes de peuplement fermé (15,4~ha)}.

\begin{figure}[H]
\centering
\includegraphics[width=0.8\linewidth]{s2_reference_window.png}
\caption{Fenêtre radar de $20\times20$~km et coupes de référence Sentinel-2.}
\label{fig-window}
\end{figure}

\subsection{Séparabilité coupe / forêt à géométrie égale}

Le \cref{tab-core} et la \cref{fig-summary} comparent les deux bandes à géométrie égale.

\begin{table}[H]
\centering\small
\caption{AUC coupe / forêt intacte [IC 95~\% par bootstrap en grappes], 7 coupes de peuplement.}
\label{tab-core}
\resizebox{\linewidth}{!}{\begin{tabular}{lcccc}
\toprule
Indicateur & L A173 & C A59 & L D66 & C D37 \\
\midrule
$\Delta\gamma^0$ pol. croisée & 0{,}80 {\scriptsize [0{,}61--0{,}93]} & 0{,}73 {\scriptsize [0{,}65--0{,}86]} & 0{,}83 {\scriptsize [0{,}66--0{,}93]} & 0{,}83 {\scriptsize [0{,}78--0{,}91]} \\
$\Delta\gamma^0$ co-pol. & 0{,}70 {\scriptsize [0{,}60--0{,}86]} & 0{,}63 {\scriptsize [0{,}53--0{,}77]} & 0{,}71 {\scriptsize [0{,}60--0{,}82]} & 0{,}74 {\scriptsize [0{,}67--0{,}82]} \\
$\gamma^0$ après, pol. croisée & 0{,}89 {\scriptsize [0{,}87--0{,}91]} & 0{,}76 {\scriptsize [0{,}72--0{,}84]} & 0{,}85 {\scriptsize [0{,}81--0{,}91]} & 0{,}82 {\scriptsize [0{,}75--0{,}89]} \\
cohérence minimale (co-pol.) & 0{,}60 {\scriptsize [0{,}53--0{,}66]} & 0{,}51 {\scriptsize [0{,}50--0{,}55]} & 0{,}81 {\scriptsize [0{,}75--0{,}86]} & 0{,}52 {\scriptsize [0{,}50--0{,}55]} \\
cohérence après (co-pol.) & 0{,}66 {\scriptsize [0{,}51--0{,}76]} & 0{,}68 {\scriptsize [0{,}61--0{,}79]} & 0{,}54 {\scriptsize [0{,}50--0{,}78]} & 0{,}74 {\scriptsize [0{,}66--0{,}83]} \\
\bottomrule
\end{tabular}
}
\end{table}

\begin{figure}[H]
\centering
\includegraphics[width=\linewidth]{m8_summary_L_vs_C.png}
\caption{AUC des cinq indicateurs clés, bande L (vert) contre bande C (violet), en montante et en descendante.}
\label{fig-summary}
\end{figure}

\begin{table}[H]
\centering\small
\caption{Moyennes coupe / forêt intacte.}
\label{tab-means}
\begin{tabular}{lcccc}
\toprule
Coupe / forêt intacte & L A173 & C A59 & L D66 & C D37 \\
\midrule
$\Delta\gamma^0$ croisée (dB) & -2{,}49 / -0{,}23 & -1{,}76 / -0{,}44 & -2{,}66 / -0{,}33 & -2{,}59 / -0{,}60 \\
cohérence min. & 0{,}28 / 0{,}32 & 0{,}20 / 0{,}21 & 0{,}34 / 0{,}48 & 0{,}16 / 0{,}16 \\
\bottomrule
\end{tabular}

\end{table}

\begin{itemize}
\item \textbf{Rétrodiffusion croisée.} La baisse est plus forte en bande L ($-2{,}5$ et $-2{,}7$~dB contre
  $-1{,}8$ et $-2{,}6$~dB en C), et le d de Cohen est plus grand (1,89 et 2,13 contre 1,03 et 1,63). L'AUC
  est meilleure en L en montante (0,80 contre 0,73) et égale en descendante (0,83). Avec 7 coupes, les IC se
  recouvrent : aucune différence d'AUC en rétrodiffusion n'est significative.
\item \textbf{Polarisation.} La polarisation croisée sépare mieux que la co-polarisation dans les quatre
  géométries (H4).
\item \textbf{Contrôle avant coupe.} Les variables de changement restent près du hasard (AUC 0,50 à 0,57).
  Les variables de niveau avant la coupe atteignent 0,54 à 0,64 (0,70 à 0,75 avec la référence
  \enquote{vérifiée}) : les parcelles coupées sont des peuplements particuliers, surtout des résineux mûrs.
  Le $\gamma^0$ après la coupe (AUC 0,85 à 0,89) est donc en partie gonflé par le type de peuplement ;
  l'indicateur principal reste la variation $\Delta\gamma^0$.
\item \textbf{Pixels de bord.} Avec des objets érodés de 20~m, toutes les AUC augmentent de 0,03 à 0,07,
  sans changer l'ordre L / C.
\end{itemize}

\subsection{Cohérence interférométrique}

\begin{table}[H]
\centering\small
\caption{Cohérence médiane (co-polarisation) de la forêt intacte et des coupes après la coupe, sur l'ensemble
des paires à 12~jours (24~jours pour une paire NISAR par géométrie).}
\label{tab-coh}
\begin{tabular}{lcc}
\toprule
Géométrie & Forêt intacte & Coupes, après la coupe \\
\midrule
L A173 & 0,35 -- 0,70 & 0,41 -- 0,74 \\
L D66 & 0,48 -- 0,71 & 0,59 -- 0,71 \\
C A59 & 0,30 -- 0,36 & 0,34 -- 0,53 \\
C D37 & 0,25 -- 0,27 & 0,27 -- 0,38 \\
\bottomrule
\end{tabular}
\end{table}

En bande C, la forêt reste au plancher de l'estimateur (0,25 à 0,36 pour une fenêtre $10\times3$) : elle est
complètement décorrélée en 12 jours, comme dans les Vosges. En bande L, elle garde 0,5 à 0,7
(\cref{tab-coh}, H1). Une coupe y devient donc un événement : en D66, la cohérence minimale des coupes
tombe à 0,34 contre 0,48 pour la forêt (AUC 0,81 contre 0,52 en C D37). La paire du 29/08 au 10/09, qui
encadre trois coupes, donne 0,45 sur les coupes et 0,63 en forêt.

En montante (A173), l'avantage est faible (AUC 0,60) : deux paires de la fenêtre de coupe (13/08 $\to$
25/08 et 25/08 $\to$ 06/09) décorrèlent toute la forêt (0,35), probablement à cause d'un épisode de pluie
ou de vent, que nous n'avons pas vérifié. Enfin, en bande C, la coupe ne se voit qu'\emph{après} : le sol nu
redevient un peu cohérent en octobre--novembre (jusqu'à 0,53), alors que la forêt reste au plancher. D'où une
AUC de la cohérence après coupe plus élevée en C (0,68--0,74) qu'en L (0,54--0,66), mais 4 à 8 semaines
après la coupe.

\subsection{La méthode du guide transposée au radar}

\begin{table}[H]
\centering\small
\caption{Méthode du guide sur $\Delta\gamma^0$ croisée (meilleur degré par géométrie) : matrice de confusion
en surface, part des coupes trouvées (au moins 25~\% de leurs pixels) et fausses alertes en part de la forêt.}
\label{tab-guide}
\resizebox{\linewidth}{!}{\begin{tabular}{lllrrrcr}
\toprule
Géométrie & Variable & Degré & Omission (\%) & Commission (\%) & Kappa & Objets trouvés & Fausses alertes (\% forêt) \\
\midrule
L A173 & d\_g0\_HV & $\geq$2 & 55 & 62 & 0{,}41 & 5/7 & 0{,}09 \\
C A59 & d\_g0\_VH & $\geq$2 & 83 & 94 & 0{,}08 & 3/7 & 0{,}32 \\
L D66 & d\_g0\_HV & $\geq$2 & 51 & 64 & 0{,}41 & 5/7 & 0{,}10 \\
C D37 & d\_g0\_VH & $\geq$3 & 84 & 83 & 0{,}16 & 2/7 & 0{,}10 \\
\bottomrule
\end{tabular}
}
\end{table}

\begin{figure}[H]
\centering
\includegraphics[width=\linewidth]{m6_guide_maps.png}
\caption{Degrés de certitude du guide (jaune 1, orange 2, rouge 3) : référence Sentinel-2, puis bande L et
bande C dans chaque direction. En haut, la fenêtre entière ; en bas, un zoom de $4\times4$~km (cadre bleu).
Contour noir : coupes de référence vérifiées.}
\label{fig-guide-maps}
\end{figure}

La bande L transpose bien mieux la méthode du guide : kappa 0,41 et 5 coupes sur 7 trouvées dans les deux
directions, contre un kappa de 0,08 à 0,16 et 2 à 3 coupes en bande C (\cref{tab-guide}, \cref{fig-guide-maps}).
Au même degré ($\geq 2$), la bande L produit 11 à 13~ha de fausses alertes sur 12\,800~ha évalués
(0,1~\% de la forêt), la bande C 41 à 52~ha. Ce résultat tient pour les quatre références (\cref{tab-refs-guide}).
On reste toutefois loin de l'objectif du guide (10~\% d'omission et de commission) : 51 à 55~\% d'omission et
62 à 64~\% de commission en L.

\begin{table}[H]
\centering\small
\caption{Kappa de la méthode du guide (meilleure combinaison) selon la référence.}
\label{tab-refs-guide}
\begin{tabular}{lcccc}
\toprule
Référence (objets) & L A173 & C A59 & L D66 & C D37 \\
\midrule
automatique (60) & 0,13 & 0,04 & 0,13 & 0,09 \\
vérifiées + probables (19) & 0,27 & 0,07 & 0,27 & 0,15 \\
vérifiées (10) & 0,34 & 0,08 & 0,34 & 0,19 \\
peuplement (7) & 0,41 & 0,08 & 0,41 & 0,16 \\
\bottomrule
\end{tabular}
\end{table}

\paragraph{Audit des fausses alertes.} Les 15 plus grands objets détectés par NISAR D66 hors de toute coupe
connue ont été photo-interprétés (\cref{fig-audit}). Deux sont des coupes tardives probables (F01, 2,5~ha ;
F04, 1,6~ha), coupées après la fenêtre \enquote{après} de Sentinel-2 ; les 13 autres sont des sénescences, des
lisières, des jaunissements de mélèze ou du bocage. Le même audit (\cref{fig-audit-s1}) sur Sentinel-1 D37 ($\Delta\gamma^0$ VH, degré 3, 15 plus grands objets) ne trouve \textbf{aucune coupe tardive} : 8 objets sur 15 (22~ha sur 36) sont des parcelles déjà ouvertes en juin (anciennes coupes, landes, prairies encore classées en forêt fermée par la BD Forêt), où la strate basse change (sénescence, travail du sol) ; les 7 autres sont des forêts intactes, des lisières ou des sénescences. La bande C réagit à la végétation basse, comme sur les coupes R019, R139 et R225 ouvertes avant juillet ; la bande L, non.

\begin{figure}[H]
\centering
\includegraphics[width=\linewidth]{m6b_commission_NISAR_D66F65_d_g0_HV.png}
\caption{Audit des 15 plus grandes fausses alertes NISAR D66 : Sentinel-2 en juillet, en novembre (couleurs
vraies) et en novembre (SWIR / PIR / rouge).}
\label{fig-audit}
\end{figure}

\begin{figure}[H]
\centering
\includegraphics[width=\linewidth]{m6b_commission_S1_T037_d_g0_VH.png}
\caption{Audit des 15 plus grandes fausses alertes Sentinel-1 D37 (mêmes vignettes).}
\label{fig-audit-s1}
\end{figure}

\subsection{Sensibilité à la référence}

\begin{table}[H]
\centering\small
\caption{AUC de $\Delta\gamma^0$ croisée selon la référence.}
\label{tab-refs}
\begin{tabular}{lcccc}
\toprule
Référence (objets) & L A173 & C A59 & L D66 & C D37 \\
\midrule
auto (60) & 0{,}56 & 0{,}51 & 0{,}56 & 0{,}53 \\
vérifiées + probables (19) & 0{,}63 & 0{,}61 & 0{,}63 & 0{,}67 \\
vérifiées (10) & 0{,}68 & 0{,}70 & 0{,}71 & 0{,}78 \\
peuplement (7) & 0{,}80 & 0{,}73 & 0{,}83 & 0{,}83 \\
\bottomrule
\end{tabular}

\end{table}

La qualité de la référence domine tout le reste (\cref{tab-refs}) : l'AUC passe de 0,51--0,56 avec la
référence automatique à 0,73--0,83 avec les 7 coupes de peuplement. Sans photo-interprétation, on aurait
conclu à tort que ni la bande L ni la bande C ne voient les coupes. L'ordre L / C, lui, dépend de la
référence : avec les références \enquote{vérifiées}, la bande C égale ou dépasse la bande L (0,78 contre 0,71
en descendante). Ces références contiennent trois parcelles déjà ouvertes en juillet (R019, R139, R225) : leur
HV en bande L était déjà bas et baisse peu (moins de 1,3~dB), alors que la bande C, sensible à la végétation
basse et aux rémanents, baisse de 1,7 à 3,9~dB. La bande L répond à la perte des troncs, la bande C à tout
changement de la strate basse. C'est pourquoi la référence principale se limite aux peuplements fermés.

\subsection{Datation radar des coupes}

\begin{table}[H]
\centering\small
\caption{Intervalle de coupe (première baisse durable de plus de 2~dB en polarisation croisée), coupes
de peuplement fermé ; R : résineux, M : mixte, F : feuillus.}
\label{tab-dating}
\begin{tabular}{lcccc}
\toprule
Coupe & L A173 & C A59 & L D66 & C D37 \\
\midrule
R013 (R, 5{,}1 ha) & 25/08--06/09 & -- & 29/08--10/09 & 25/08--06/09 \\
R060 (M, 2{,}9 ha) & 25/08--06/09 & 09/07--21/07 & 29/08--10/09 & 25/08--06/09 \\
R100 (R, 2{,}1 ha) & -- & 07/09--19/09 & -- & 06/09--18/09 \\
R144 (R, 1{,}6 ha) & -- & -- & -- & 25/08--06/09 \\
R161 (F, 1{,}4 ha) & 06/09--18/09 & 07/09--19/09 & 29/08--10/09 & 06/09--18/09 \\
R198 (M, 1{,}3 ha) & 20/07--13/08 & 21/07--02/08 & 24/07--17/08 & 25/08--06/09 \\
R246 (R, 1{,}1 ha) & 20/07--13/08 & 14/08--26/08 & 12/07--24/07 & 25/08--06/09 \\
\bottomrule
\end{tabular}

\end{table}

\begin{figure}[H]
\centering
\includegraphics[width=\linewidth]{m7_cut_timeseries.png}
\caption{Séries de $\gamma^0$ en polarisation croisée des coupes confirmées et médiane de la forêt intacte.}
\label{fig-ts}
\end{figure}

En bande L, les séries sont lisses et les deux géométries concordent (\cref{tab-dating}, \cref{fig-ts}) :
R013 et R060 fin août -- début septembre, R161 début septembre, R198 et R246 fin juillet -- mi-août. En bande
C, l'orbite D37 subit des sauts synchrones de $\pm3$~dB sur toute la forêt, probablement liés à l'humidité ;
elle place 5 coupes sur 7 dans la même paire, ce qui n'est pas crédible. L'orbite A59, plus stable, resserre
la date de R198 (21/07 -- 02/08). R100 n'est détectée qu'en bande C (septembre, dans les deux orbites) ; en L,
la baisse reste sous 2~dB. Une éclaircie ou un débardage tardif, auquel la bande C réagit davantage, est
l'hypothèse la plus simple.

\subsection{Bilan des hypothèses}

\begin{table}[H]
\centering\small
\caption{Bilan des hypothèses sur le Morvan.}
\label{tab-hyp}
\begin{tabular}{lp{10.5cm}l}
\toprule
H & Constat & Verdict \\
\midrule
H1 & Forêt intacte : cohérence 0,5--0,7 en L contre 0,25--0,36 en C (plancher). & confirmée \\
H2 & D66 : AUC 0,81 contre 0,52 en C ; A173 : 0,60 (décorrélation générale fin août). & partielle \\
H3 & Baisse plus forte en L (d de Cohen 1,9--2,1 contre 1,0--1,6) ; AUC L $\geq$ C, non significatif. & partielle \\
H4 & Polarisation croisée $>$ co-polarisation dans les deux bandes et les quatre géométries. & confirmée \\
\bottomrule
\end{tabular}
\end{table}

%=====================================================================
\section{Discussion}

\paragraph{Comparaison avec les Vosges.} Les deux sites donnent le même message. La forêt reste cohérente
en bande L (0,64--0,70 dans les Vosges, 0,5--0,7 ici) et pas en bande C (0,28 et 0,25--0,36). L'avantage de
la bande L en rétrodiffusion dépend de la géométrie : dans les Vosges, il n'existait qu'en montante ; ici,
l'AUC est meilleure en montante et égale en descendante. La réplication lève les deux limites du premier
test (une seule coupe, une seule orbite Sentinel-1) et confirme qu'il ne s'agissait pas d'un cas
particulier. En revanche, l'AUC de cohérence presque parfaite des Vosges (0,995) ne se retrouve pas : sur
7 coupes de dates différentes, les paires d'été sont parfois décorrélées sur toute la forêt.

\paragraph{Pourquoi la méthode du guide favorise la bande L.} À AUC égale (0,83 en descendante), la bande L
donne un kappa de 0,41 et la bande C de 0,16. L'AUC mesure le classement moyen des pixels ; une méthode à
seuils $m-k\sigma$ ne retient que la queue de la distribution, où compte l'écart exprimé en écarts-types
(d de Cohen 2,13 contre 1,63). La bande L est aussi plus stable d'une date à l'autre, ce qui resserre la
distribution de la forêt et réduit les fausses alertes. Pour un service opérationnel fondé sur des seuils,
c'est ce critère qui compte.

\paragraph{La référence, facteur limitant.} Le guide recommande le contrôle par photo-interprétation ; ce
travail montre qu'il est indispensable. En intra-saison, la sénescence des feuillus produit la plupart des
baisses de NDVI : seules 10 des 60 coupes qualifiées de la fenêtre sont des coupes certaines. Les scores
radar mesurés contre une référence automatique sont des bornes basses très pessimistes.

\paragraph{Complémentarité des deux bandes.} La bande C réagit à des événements que la bande L voit peu
(R100, sans doute une éclaircie ou un débardage) et sa revisite permet parfois de resserrer une date
(R198). La bande L donne des séries plus stables, des dates cohérentes entre géométries et moins de fausses
alertes. Un détecteur opérationnel gagnerait à combiner les deux, avec une confirmation temporelle.

\paragraph{Limites.} Sept coupes et 15~ha seulement : les IC sont larges et les différences d'AUC en
rétrodiffusion ne sont pas significatives. Une seule saison, un seul paysage dominé par les coupes de
résineux. La référence optique rate les coupes tardives (deux d'entre elles apparaissent comme fausses
alertes radar) et les coupes sans sol nu. Les estimateurs de cohérence diffèrent (NISAR GUNW contre SNAP
$10\times3$), les produits NISAR sont provisoires et la météo n'a pas été vérifiée.

%=====================================================================
\section{Conclusion}

Sur sept coupes rases vérifiées du Morvan et à géométrie égale dans les deux directions de visée :
\begin{enumerate}
\item la bande L garde la forêt cohérente (0,5--0,7) et la bande C non (plancher à 0,25--0,36) : H1 est
  confirmée sur un deuxième site ;
\item la cohérence L sépare les coupes en descendante (AUC 0,81 contre 0,52), moins en montante, où un
  épisode de décorrélation générale masque le signal ;
\item en rétrodiffusion, les AUC sont proches (L $\geq$ C), mais le contraste est plus fort en L, et la
  méthode du guide transposée au radar fonctionne bien mieux en bande L (kappa 0,41, 5 coupes sur 7) qu'en
  bande C (0,08--0,16, 2 à 3 coupes), sans atteindre l'objectif de 10~\% d'erreur ;
\item la qualité de la référence optique domine les résultats : la photo-interprétation est indispensable.
\end{enumerate}
La chaîne (paquet \texttt{nisar\_s1}, configuration par site) est directement transposable à ROSE-L~\cite{esa2020rosel}
et aux produits ISCE3~\cite{isce3}. La suite logique est un détecteur combinant cohérence L, rétrodiffusion L
et revisite C, évalué sur plusieurs saisons et plusieurs régions.
